\documentclass[10pt,a4paper]{amsart}
\usepackage[margin=19mm]{geometry}
\usepackage{amsmath,amssymb,mathtools}
\usepackage[scr=rsfs,cal=euler]{mathalfa}
\usepackage{xfrac}
\usepackage[unicode,hidelinks,bookmarks=false]{hyperref}
\newcommand{\ie}{i.e.\ }
\newcommand{\cf}{cf.\ }
\newcommand{\resp}{respectively\ }
\newcommand{\Subsection}{\subsection}
\providecommand{\nobreakdash}{\nobreak\hbox{-}}
\setlength{\emergencystretch}{2em}
\multlinegap0pt
\usepackage{amssymb}
\usepackage{mathtools}
\usepackage{xcolor}
\usepackage{bm}
\usepackage{cancel}
\usepackage[normalem]{ulem}
\usepackage{enumerate}
\newtheorem{proposition}{Proposition}
\newtheorem{lemma}{Lemma}
\newcommand{\bK}{\mathbf{K}}
\newcommand{\bL}{\mathbf{L}}
\newcommand{\del}{\partial}
\newcommand{\lra}{\longrightarrow}
\providecommand{\todo}[1]{\P \textbf{#1} \P}
\newcommand{\ve}{\varepsilon}
\title{Concavity of spacetimes}
\author{Tobias Beran, Darius Er\"os, Shin-ichi Ohta, Felix Rott}
\date{Source: arXiv:2509.26196v2 (CC BY 4.0)}
\begin{document}
\maketitle

\noindent\emph{Source and licence.} Concavity of spacetimes. Version arXiv:2509.26196v2 (CC BY 4.0), distributed by arXiv under the Creative Commons Attribution 4.0 International licence, http://creativecommons.org/licenses/by/4.0/, as recorded in the arXiv licence metadata for this exact version and shown on the abstract page https://arxiv.org/abs/2509.26196v2. Full licence notice: \texttt{licenses/arxiv-2509.26196-CC-BY-4.0.txt}.\par\smallskip

\noindent\textit{Editorial scope.} Counted unit: the article's lemma VtimeGen (source0.tex lines 897--904). The article's complete original proof of it is delivered with it (source0.tex lines 906--932, one \texttt{proof} environment of 27 lines). No mathematics is written, repaired or re-derived: the statement and the proof are the article's own lines, in the article's own notation, and no retained line is substituted or re-flowed.  Retained uncounted as the source-local context the proof needs: lines 818--824; lines 874--880.  Nothing was omitted from any retained span, so the package carries no omission record.  No retained line contains a citation, so the wrapper carries no bibliography and no bibliography entry.  No retained line contains a figure environment, an image-inclusion command or drawing code, so no image is delivered and the package declares no asset dependency.  Editorial formatting only, in addition: an AMS \texttt{amsart} preamble that restates the article's own declarations and macros used by the retained lines, namely the environments \texttt{proposition}, \texttt{lemma} on separate counters, together with the standard packages those lines need.  The retained environments print in this document as Proposition 1, Lemma 1 in order of appearance, whereas the article prints the same items with the numbering of its own full document.  The complete native source of the article as distributed in the arXiv e-print for this exact version is bundled unchanged as \texttt{sources/186068/source0.tex}.
\begin{proposition}[Concave norm]
\label{prop: concave length}
Let $(M,L)$ be a Berwald spacetime
of nonnegative flag curvature $\bK \ge 0$ in timelike directions. 
If $\sigma \colon [0,1] \times (-\ve,\ve) \lra U$ is a variation in a convex normal neighborhood $U$ such that $\sigma_s(t):=\sigma(t,s)$ is a geodesic for each $s \in (-\ve,\ve)$ and that $V(t,s):=\partial_s\sigma(t,s)$ is timelike for all $(t,s)$, then, for every $s$, the function $t\longmapsto F(V(t,s))$ is concave. 
%Then the length functional is concave along geodesic variations, i.e., if $\sigma : [0,1] \times [-\varepsilon, \varepsilon] \to M$ is a geodesic variation with a fixed starting point, then $\bL'' \leq 0$.
\end{proposition}

\begin{lemma}[Timelike homotopy]
\label{lem: timelike homotopy}
Let $U$ be a convex normal neighborhood in a Finsler spacetime $(M,L)$.
Then, for any timelike curves $\alpha,\beta \colon [0,1] \lra U$,
there exists a homotopy $h\colon [0,1] \times [0,1] \lra U$ such that $h(0;s)=\alpha(s)$, $h(1;s)=\beta(s)$ and $s \longmapsto h(t;s)$ is timelike for each $t$,
and that $h$ and $\partial_s h$ are continuous.
\end{lemma}

\begin{lemma}[Timelike variation]
\label{lem: VtimeGen}
Let $(M,L)$ be a Berwald spacetime of nonnegative flag curvature $\bK \ge 0$ in timelike directions,
and $U \subset M$ be a convex normal neighborhood.
Given timelike curves $\alpha,\beta\colon [0,1] \lra U$, let $t\longmapsto \sigma(t,s)$ be the geodesic from $\alpha(s)$ to $\beta(s)$ in $U$.
Then, $\partial_s \sigma$ is timelike and $t \longmapsto F(\del_s \sigma(t,s))$ is concave for each $s \in [0,1]$.
%We can also choose one of $\alpha$ and $\beta$ to be constant, and then $\partial_s \sigma$ is timelike except for $t=0$ or $t=1$, respectively.
\end{lemma}

\begin{proof}
Let $h(\lambda;s)$ be a homotopy between $\alpha$ and $\beta$ given by Lemma~\ref{lem: timelike homotopy}, i.e.,
$h(0;s)=\alpha(s)$, $h(1;s)=\beta(s)$ and $\partial_s h$ is always timelike.
For each $\lambda \in [0,1]$, we consider the variation $\Xi(\lambda;\cdot,\cdot) \colon [0,1] \times [0,1] \lra U$ such that $t \longmapsto \Xi(\lambda;t,s)$ is the geodesic from $\Xi(\lambda;0,s):=h(\lambda;s)$ to $\Xi(\lambda;1,s):=\beta(s)$ in $U$.
Observe that $\Xi(0;t,s)=\sigma(t,s)$ and $\Xi(1;t,s)= \beta(s)$.

Now, we set
\[ \lambda_0:=\sup\bigl\{ \lambda \in [0,1] \mid \del_s \Xi(\lambda;t,s) \text{ is not timelike for some } t,s \in [0,1] \bigr\}. \]
Note that $\lambda_0 < 1$, since $\partial_s\Xi(1;t,s)=\dot\beta(s)$ is timelike and being timelike is an open condition.
%Moreover, we can find $s_0,t_0$ such that $\del_s \Xi(r_0;s_0,t_0)$ is not timelike and thus null. \todo{can now drop this sentence}
%
Then, for any $\lambda \in (\lambda_0,1]$, since $\del_s \Xi(\lambda;t,s)$ is timelike for all $t,s \in [0,1]$, it follows from Proposition~\ref{prop: concave length} that
$F(\del_s \Xi(\lambda;t,s))$ is concave in $t$.
Taking the limit as $\lambda \to \lambda_0$, we find that $F(\del_s \Xi(\lambda_0;t,s))$ is concave in $t$.
Hence,
\begin{align*}
F\bigl( \del_s \Xi(\lambda_0;t,s) \bigr)
&\ge (1-t)F\bigl( \del_s \Xi(\lambda_0;0,s) \bigr) +tF\bigl( \del_s \Xi(\lambda_0;1,s)\bigr) \\
&= (1-t)F\bigl( \del_s h(\lambda_0;s)\bigr) +tF\bigl( \dot{\beta}(s) \bigr) >0
\end{align*}
for all $t,s \in [0,1]$.
This shows that $\partial_s\Xi(\lambda_0;t,s)$ is also timelike for all $t,s \in [0,1]$.

If $\lambda_0>0$, then we deduce from the openness of being a timelike vector that $\partial_s\Xi(\lambda;t,s)$ is timelike for all $t,s \in [0,1]$ and $\lambda \in (\lambda_0-\ve,1]$ for sufficiently small $\ve>0$, contradicting the definition of $\lambda_0$.
Consequently, $\partial_s\Xi(\lambda;t,s)$ is timelike for all $\lambda,t,s \in [0,1]$.
Taking $\lambda=0$ implies that $\del_s \sigma$ is timelike, and $F(\del_s \sigma(t,s))=F(\del_s \Xi(0;t,s))$ is concave in $t$. 
\end{proof}

\end{document}
